% ===================================================================
%  図表第 11 号 — 町とその建物。— 火事。
%
%  Ceci n'est PAS une transcription : c'est une traduction, faite en
%  2026, en japonais d'aujourd'hui. Voir texto/ja/10-zuhyo-01.tex
%  pour la consigne, qui vaut pour les seize fichiers.
% ===================================================================

%%K t11-apar-1 apar
\VUsaut{2.0mm}
\VUcentre{13.2pt}{90}{第三部}
\VUsaut{3.0mm}
\VUfilet{16mm}
\VUsaut{5.6mm}
\VUtitre{15.0pt}{60}{図表第 11 号}
\VUsaut{1.4mm}
\VUpk{11.4pt}{40}{町とその建物。--- 火事。}
\VUsaut{2.2mm}

%%K t11-tit-1 sub
\VUpk{10.2pt}{40}{町はずれ。}
\VUsaut{3.0mm}

%%K t11-c1-01-1 p
1. --- 私は大きな町に住んでいる。海から百キロメートル離れているのに、それでも
第一等の\VUgras{港} \textsuperscript{(1)} である。町のかたわらを流れる川は幅四百
メートルあり、よい泊地をなしている。

%%K t11-c1-02-1 p
2. --- \VUgras{波止場} \textsuperscript{(2)} に沿った一帯はまったく海港と工場の
姿をしており、\VUgras{町はずれ} \textsuperscript{(3)} をなしていて、住むのは水夫と
職人ばかりである。職人は\VUgras{工場} \textsuperscript{(4)} と\VUgras{ガス場}
\textsuperscript{(5)} で働く。ガス場の高い\VUgras{煙突} \textsuperscript{(6)} と
\VUgras{ガス槽}は遠くからでも見える。

%%K t11-c1-03-1 p
3. --- 中世には町は城壁に囲まれていて、いまもその名残が幾らか見つかる。とくに
あの\VUgras{門} \textsuperscript{(8)} で、古い\VUgras{城砦} \textsuperscript{(7)} の門
である。その両側に\VUgras{塔} \textsuperscript{(9)} が二つあり、いまは\VUgras{牢}に
使われている。

%%K t11-c1-04-1 p
4. --- この\VUgras{界隈}はすべて古びて見えるが、一つだけ比較的新しい建物が
ある。\VUgras{税関} \textsuperscript{(10)} である。波止場とあの小さな\VUgras{通}
\textsuperscript{(11)} とのあいだにあり、その通は古い\VUgras{町役場}
\textsuperscript{(12)} へ通じている。その役場はたやすく見分けられる。旧市街の上に
そびえる巨大な\VUgras{鐘楼} \textsuperscript{(13)} があるからである。

%%K t11-c1-05-1 p
5. --- 通の向い側に税関と同じ様式の建物が\VUgras{立って}いる。\VUgras{取引所}
\textsuperscript{(14)} である。その一面は小さな広場に向いていて、そこに
\VUgras{県庁} \textsuperscript{(15)} がある。我々の町は都ではないが、それでも大切な
郡都である。

%%K t11-tit-2 sub
\VUsaut{5.0mm}
\VUpk{10.2pt}{40}{町の高い方。}
\VUsaut{3.4mm}

%%K t11-c1-06-1 p
6. --- 守備の兵は数が少なくなく、広くたいそう衛生のよい\VUgras{兵舎}
\textsuperscript{(16)} に住んでいる。兵舎は\VUgras{町の公園} \textsuperscript{(17)} の
柵まで続いている。この公園へ通じる\VUgras{並木の大通} \textsuperscript{(19)} は
\VUgras{貯蓄銀行} \textsuperscript{(18)} の裏を過ぎ、\VUgras{大聖堂}
\textsuperscript{(20)} の広場に出る。

%%K t11-c1-07-1 p
7. --- これらの建物はみな小さな丘に段をなして建っており、我々の大きな
\VUgras{中学校} \textsuperscript{(21)} もそこにある。

%%K t11-c1-07-2 p
少し傾いた道を下って大通に出ると、\VUgras{裁判所} \textsuperscript{(22)} に着く。
その前には心地よい\VUgras{公けの庭} \textsuperscript{(26)} が広がっている。この
\VUgras{広場}は大きくはないが、町の真中に緑と涼しさの沃地をなしている。だから
町の人がよく来る。\VUgras{カフェ} \textsuperscript{(25)} の日除の下に坐って、木曜と
日曜に芝と花壇のあいだの\VUgras{奏楽堂} \textsuperscript{(27)} で演じる軍楽を聴くのが
好きなのである。

%%K t11-c1-08-1 p
8. --- その芝の一つに青銅の\VUgras{像} \textsuperscript{(28)} が立っている。我々の
町の人で、生れた町に大きな功のあった者を記念して建てられたものである。

%%K t11-tit-3 sub
\VUsaut{4.0mm}
\VUpk{10.2pt}{40}{行き来。}
\VUsaut{3.4mm}

%%K t11-c1-09-1 p
9. --- 我々の町にはまだ電の灯がなく、\VUgras{ガス燈} \textsuperscript{(29)} しか
ない。しかし\VUgras{土地の新聞}はこの明りのしかけを改めよと唱えている。
\VUgras{ちょうど}電車の\VUgras{牽きかた}がすでに改まったようにである。
\VUgras{あの}\VUgras{馬}に\VUgras{曳かれた}古い車は、便利な\VUgras{電車}
\textsuperscript{(31)} に替わった。たいそう\VUgras{安い}賃で、客を\VUgras{町の
一方から他方へ}早く運ぶ。

%%K t11-c1-10-1 p
10. --- 裁判所のそばが\VUgras{銀行} \textsuperscript{(32)} で、商いの手形をそこで
割り引く。\VUgras{銀行の集金人} \textsuperscript{(33)} が家々を回って取り立てに
行く。

%%K t11-tit-4 sub
\VUsaut{4.0mm}
\VUpk{10.2pt}{40}{美術と学問。}
\VUsaut{3.4mm}

%%K t11-c1-11-1 p
11. --- そのそばの美しい建物は\VUgras{博物館}である。中には町の\VUgras{図書館}
\textsuperscript{(34)} と、美しい\VUgras{博物の標本} \textsuperscript{(35)}（自然史の
館）と、洒落た\VUgras{画の間} \textsuperscript{(36)} とが一緒にあり、それは見事な
常設の\VUgras{展覧}をなしている。

%%K t11-c1-11-2 p
週に二度、人々に開かれるが、見物人は\VUgras{番人} \textsuperscript{(37)} に少しの
心づけをやれば、どの日にも見られる。\VUgras{画家} \textsuperscript{(38)} たちが
そこへ描きに来る。画架の\VUgras{前に}立ち、手に\VUgras{絵皿}を持って、名高い
絵を写しているのが見える。

%%K t11-c1-12-1 p
12. --- 博物館の裏の高みには、あの\VUgras{円屋根} \textsuperscript{(40)} --- それは
\VUgras{病院} \textsuperscript{(39)} の円屋根である --- と、\VUgras{師範学校}
\textsuperscript{(41)} の建物とがかろうじて見える。我々の町の小学校の先生はそこで
仕込まれたのである。芝居小屋の広場には\VUgras{郵便局} \textsuperscript{(43)} が
あり、\VUgras{私の家} \textsuperscript{(42)} に置かれている。

%%K t11-tit-5 sub
\VUsaut{5.0mm}
\VUpk{11.4pt}{40}{火事。}
\VUsaut{3.0mm}

%%K t11-tit-6 sub
\VUpk{10.2pt}{40}{「火事だ」の声。}
\VUsaut{3.4mm}

%%K t11-c2-01-1 p
1. --- 恐ろしい知らせが町に広がる。芝居小屋にいま火が出たというのである。私が
\VUgras{広場} \textsuperscript{(96)} に着いたときには、人はもう\VUgras{歩道}
\textsuperscript{(46)} と\VUgras{車道} \textsuperscript{(47)} に詰めかけていた。
\VUgras{郵便の役人} \textsuperscript{(44)} と\VUgras{郵便配達} \textsuperscript{(45)}
が郵便局から出てくる。一人の\VUgras{電車の運転手} \textsuperscript{(48)} は車を
止めねばならなかった。町の\VUgras{病人運びの車} \textsuperscript{(50)} を通すため
である。その車は\VUgras{外科医} \textsuperscript{(52)} と\VUgras{看護婦}
\textsuperscript{(51)} を乗せて来ている。一人の\VUgras{怪我人} \textsuperscript{(53)} を
手当するためで、その人は\VUgras{担架} \textsuperscript{(54)} で\VUgras{新聞の売店}
\textsuperscript{(55)} の近くまで運ばれてきた。

%%K t11-c2-02-1 p
2. --- 調練から戻ってきた歩兵の一隊が足を止め、馬に乗った\VUgras{町の警手}
\textsuperscript{(61)} とともに秩序を保つのを助けている。\VUgras{騎馬の警手}の馬は
立ち上がる。彼らは\VUgras{兵} \textsuperscript{(57)} や\VUgras{憲兵}
\textsuperscript{(63)} よりもうまく\VUgras{見物人} \textsuperscript{(56)} を退かせる
ことができる。憲兵はそれに、たいそう忙しい。その一人が\VUgras{巡査}
\textsuperscript{(64)} に、もう一人の\VUgras{怪我人} \textsuperscript{(65)} をどこへ
運ぶかを告げている。梯子から落ちた勇ましい消防夫である。

%%K t11-c2-03-1 p
3. --- \VUgras{士官} \textsuperscript{(66)} が、いま着く\VUgras{蒸汽喞筒}
\textsuperscript{(67)} をどこに据えるかを示している。その力強い機械を待つあいだ、
\VUgras{手押喞筒} \textsuperscript{(68)} を据えさせた。その長い\VUgras{水管}
\textsuperscript{(69)} が火の上に水を注いでいる。六人の消防夫がそれを押し、その
仲間が\VUgras{筒先} \textsuperscript{(70)} を握って\VUgras{水の柱}
\textsuperscript{(71)} を火の真中へ向けている。

%%K t11-c2-04-1 p
4. --- 芝居小屋の反対側に大きな\VUgras{梯子} \textsuperscript{(72)} が立てられた。
それがあってこそ消防夫は、黒い\VUgras{煙} \textsuperscript{(75)} の渦のあいだから
噴き上がる炎に近づけるのである。

%%K t11-tit-7 sub
\VUsaut{5.0mm}
\VUpk{10.2pt}{40}{勇ましい振舞。}
\VUsaut{3.4mm}

%%K t11-c2-05-1 p
5. --- \VUgras{火} \textsuperscript{(74)} は\VUgras{芝居小屋} \textsuperscript{(73)}
の玄関で出た。ちょうど\VUgras{歌劇}の稽古の最中で、その初めての\VUgras{興行}が
明日のはずであった。幸い場内には\VUgras{見物} \textsuperscript{(76)} が数人しか
いなかった。この気の毒な人々は\VUgras{テラス} \textsuperscript{(77)} に逃れており、
勇ましい\VUgras{消防夫} \textsuperscript{(86)} が梯子と綱を持って助けに来ている。

%%K t11-c2-06-1 p
6. --- \VUgras{車寄せ} \textsuperscript{(78)} の下には興行を知らせる\VUgras{ちらし}
\textsuperscript{(79)} が貼ってある。そこには\VUgras{楽人} \textsuperscript{(81)} が
楽器を抱えて逃げ出すのが見える。\VUgras{楽隊} \textsuperscript{(80)} の楽人である。
\VUgras{提琴} \textsuperscript{(81)}、\VUgras{横笛} \textsuperscript{(82)}、
\VUgras{大提琴} \textsuperscript{(83)}、\VUgras{喇叭} \textsuperscript{(84)}、
\VUgras{太鼓} \textsuperscript{(85)} などである。\VUgras{道具} \textsuperscript{(87)} の
一部を助け出そうと人々が努めている。

%%K t11-c2-07-1 p
7. --- \VUgras{男役者} \textsuperscript{(89)} と\VUgras{女役者} \textsuperscript{(88)}、
男女の\VUgras{歌い手}は、舞台の\VUgras{衣裳} \textsuperscript{(90)} のまま逃げ出す
ほかなかった。第一の歌い手が\VUgras{新聞記者} \textsuperscript{(92)} に、事の起こり
を語っている。\VUgras{その取材の人}は初めから役所の人々とともに駆けつけて
いた。\VUgras{知事} \textsuperscript{(91)}、\VUgras{町長} \textsuperscript{(93)}、
\VUgras{警察署長} \textsuperscript{(94)}、それに\VUgras{判事} \textsuperscript{(95)} で
ある。判事は誰に責めがあるかを調べることになる。
\VUsaut{6.0mm}
\VUfilet{22mm}
